% =====================================================================================
%  Big Data (ULPGC, 2026-2027) -- Individual Assignment 03
%  Author: Daniel Rodríguez Alonso
%  Build:  cd report && latexmk -pdf main.tex      (or: make report from the folder root)
%  Every number in this file comes from report/generated/*.tex (written by analysis/make_all.py)
%  or is traceable to results/ (raw CSVs, evidence files).
% =====================================================================================
\documentclass[conference,a4paper]{IEEEtran}
\IEEEoverridecommandlockouts

\usepackage[utf8]{inputenc}
\usepackage[T1]{fontenc}
\usepackage{amsmath,amssymb}
\usepackage{graphicx}
\usepackage{booktabs}
\usepackage{siunitx}
\sisetup{group-separator={\,}, group-minimum-digits=4}
\usepackage[caption=false,font=footnotesize]{subfig}
\usepackage{xcolor}
\usepackage{listings}
\usepackage{microtype}
\usepackage{url}
\usepackage[hidelinks]{hyperref}
\usepackage[capitalise,noabbrev]{cleveref}

\lstset{basicstyle=\ttfamily\footnotesize, columns=fullflexible, breaklines=true,
        frame=single, numbers=none, keepspaces=true}
\graphicspath{{figures/}}

% Key numbers generated by the analysis scripts (\newcommand{\Something}{value})
% Generated by analysis/make_all.py -- do not edit by hand.
\newcommand{\NumRowsTotal}{3480}  % all timing rows incl. warm-ups
\newcommand{\NumRowsMeasured}{3210}  % measured ok rows of the timing experiments
\newcommand{\NumWarmupRows}{270}
\newcommand{\NumFailedRows}{0}
\newcommand{\NumThreadMismatch}{0}  % rows where the real team size != requested p
\newcommand{\ChecksumDistinctMax}{1}  % distinct checksums per n across seq/omp/scalar/avx2 builds and all p
\newcommand{\NumProcessRuns}{480}
\newcommand{\MachineMinutes}{22}  % sum of process wall times of all experiments
\newcommand{\OtherBusyMedian}{0.06}
\newcommand{\OtherBusyPninety}{1.27}
\newcommand{\NumContaminated}{48}
\newcommand{\NumLongReps}{414}
\newcommand{\ContaminationThreshold}{1.0}
\newcommand{\LoadMax}{9.0}
\newcommand{\SensMaxRelDiffPct}{3.0}  % max |relative change| of S(p) when contaminated reps are excluded
\newcommand{\NLargest}{3072}
\newcommand{\SLargeOne}{0.82}
\newcommand{\SLargeLoOne}{0.81}
\newcommand{\SLargeHiOne}{0.85}
\newcommand{\ELargeOne}{0.82}
\newcommand{\TLargeOne}{16.7}
\newcommand{\TLargeIqrOne}{0.0961}
\newcommand{\SLargeTwo}{1.57}
\newcommand{\SLargeLoTwo}{0.90}
\newcommand{\SLargeHiTwo}{1.63}
\newcommand{\ELargeTwo}{0.79}
\newcommand{\TLargeTwo}{8.73}
\newcommand{\TLargeIqrTwo}{0.239}
\newcommand{\SLargeThree}{2.14}
\newcommand{\SLargeLoThree}{1.67}
\newcommand{\SLargeHiThree}{2.21}
\newcommand{\ELargeThree}{0.71}
\newcommand{\TLargeThree}{6.41}
\newcommand{\TLargeIqrThree}{0.334}
\newcommand{\SLargeFour}{2.69}
\newcommand{\SLargeLoFour}{2.43}
\newcommand{\SLargeHiFour}{2.87}
\newcommand{\ELargeFour}{0.67}
\newcommand{\TLargeFour}{5.10}
\newcommand{\TLargeIqrFour}{0.517}
\newcommand{\SLargeFive}{3.18}
\newcommand{\SLargeLoFive}{2.34}
\newcommand{\SLargeHiFive}{3.49}
\newcommand{\ELargeFive}{0.64}
\newcommand{\TLargeFive}{4.31}
\newcommand{\TLargeIqrFive}{0.332}
\newcommand{\SLargeSix}{3.97}
\newcommand{\SLargeLoSix}{2.55}
\newcommand{\SLargeHiSix}{4.10}
\newcommand{\ELargeSix}{0.66}
\newcommand{\TLargeSix}{3.46}
\newcommand{\TLargeIqrSix}{1.39}
\newcommand{\SLargeEight}{2.97}
\newcommand{\SLargeLoEight}{2.29}
\newcommand{\SLargeHiEight}{3.07}
\newcommand{\ELargeEight}{0.37}
\newcommand{\TLargeEight}{4.63}
\newcommand{\TLargeIqrEight}{0.391}
\newcommand{\SLargeTen}{3.40}
\newcommand{\SLargeLoTen}{2.27}
\newcommand{\SLargeHiTen}{3.61}
\newcommand{\ELargeTen}{0.34}
\newcommand{\TLargeTen}{4.04}
\newcommand{\TLargeIqrTen}{1.75}
\newcommand{\SLargeTwelve}{4.11}
\newcommand{\SLargeLoTwelve}{3.66}
\newcommand{\SLargeHiTwelve}{4.25}
\newcommand{\ELargeTwelve}{0.34}
\newcommand{\TLargeTwelve}{3.35}
\newcommand{\TLargeIqrTwelve}{0.257}
\newcommand{\SLargeSixteen}{3.51}
\newcommand{\SLargeLoSixteen}{2.86}
\newcommand{\SLargeHiSixteen}{3.75}
\newcommand{\ELargeSixteen}{0.22}
\newcommand{\TLargeSixteen}{3.91}
\newcommand{\TLargeIqrSixteen}{0.492}
\newcommand{\SLargeTwentyFour}{3.93}
\newcommand{\SLargeLoTwentyFour}{3.47}
\newcommand{\SLargeHiTwentyFour}{4.13}
\newcommand{\ELargeTwentyFour}{0.16}
\newcommand{\TLargeTwentyFour}{3.49}
\newcommand{\TLargeIqrTwentyFour}{0.0994}
\newcommand{\TseqLarge}{13.7}
\newcommand{\TseqLargeIqr}{0.144}
\newcommand{\KseqLarge}{5}
\newcommand{\KpLarge}{5}
\newcommand{\PBestLarge}{12}
\newcommand{\SBestLarge}{4.11}
\newcommand{\STwoNTwoFiftySix}{1.56}
\newcommand{\ETwoNTwoFiftySix}{0.78}
\newcommand{\SSixNTwoFiftySix}{4.22}
\newcommand{\ESixNTwoFiftySix}{0.70}
\newcommand{\STwelveNTwoFiftySix}{4.36}
\newcommand{\ETwelveNTwoFiftySix}{0.36}
\newcommand{\STwentyFourNTwoFiftySix}{4.08}
\newcommand{\ETwentyFourNTwoFiftySix}{0.17}
\newcommand{\SBestNTwoFiftySix}{4.36}
\newcommand{\PBestNTwoFiftySix}{12}
\newcommand{\TseqNTwoFiftySix}{0.0079}
\newcommand{\STwoNFiveTwelve}{1.59}
\newcommand{\ETwoNFiveTwelve}{0.79}
\newcommand{\SSixNFiveTwelve}{4.12}
\newcommand{\ESixNFiveTwelve}{0.69}
\newcommand{\STwelveNFiveTwelve}{4.25}
\newcommand{\ETwelveNFiveTwelve}{0.35}
\newcommand{\STwentyFourNFiveTwelve}{4.43}
\newcommand{\ETwentyFourNFiveTwelve}{0.18}
\newcommand{\SBestNFiveTwelve}{4.43}
\newcommand{\PBestNFiveTwelve}{24}
\newcommand{\TseqNFiveTwelve}{0.0642}
\newcommand{\STwoNTenTwentyFour}{1.58}
\newcommand{\ETwoNTenTwentyFour}{0.79}
\newcommand{\SSixNTenTwentyFour}{3.88}
\newcommand{\ESixNTenTwentyFour}{0.65}
\newcommand{\STwelveNTenTwentyFour}{4.06}
\newcommand{\ETwelveNTenTwentyFour}{0.34}
\newcommand{\STwentyFourNTenTwentyFour}{4.25}
\newcommand{\ETwentyFourNTenTwentyFour}{0.18}
\newcommand{\SBestNTenTwentyFour}{4.25}
\newcommand{\PBestNTenTwentyFour}{24}
\newcommand{\TseqNTenTwentyFour}{0.505}
\newcommand{\STwoNTwentyFortyEight}{1.59}
\newcommand{\ETwoNTwentyFortyEight}{0.80}
\newcommand{\SSixNTwentyFortyEight}{3.80}
\newcommand{\ESixNTwentyFortyEight}{0.63}
\newcommand{\STwelveNTwentyFortyEight}{4.05}
\newcommand{\ETwelveNTwentyFortyEight}{0.34}
\newcommand{\STwentyFourNTwentyFortyEight}{4.05}
\newcommand{\ETwentyFourNTwentyFortyEight}{0.17}
\newcommand{\SBestNTwentyFortyEight}{4.05}
\newcommand{\PBestNTwentyFortyEight}{24}
\newcommand{\TseqNTwentyFortyEight}{4.07}
\newcommand{\STwoNThirtySeventyTwo}{1.57}
\newcommand{\ETwoNThirtySeventyTwo}{0.79}
\newcommand{\SSixNThirtySeventyTwo}{3.97}
\newcommand{\ESixNThirtySeventyTwo}{0.66}
\newcommand{\STwelveNThirtySeventyTwo}{4.11}
\newcommand{\ETwelveNThirtySeventyTwo}{0.34}
\newcommand{\STwentyFourNThirtySeventyTwo}{3.93}
\newcommand{\ETwentyFourNThirtySeventyTwo}{0.16}
\newcommand{\SBestNThirtySeventyTwo}{4.11}
\newcommand{\PBestNThirtySeventyTwo}{12}
\newcommand{\TseqNThirtySeventyTwo}{13.7}
\newcommand{\SIkjSixNTenTwentyFour}{5.85}
\newcommand{\SIkjTwelveNTenTwentyFour}{5.03}
\newcommand{\SIkjBestNTenTwentyFour}{5.85}
\newcommand{\PIkjBestNTenTwentyFour}{6}
\newcommand{\GflopsIkjBestNTenTwentyFour}{19.5}
\newcommand{\TseqIkjNTenTwentyFour}{0.644}
\newcommand{\GflopsBlockedBestNTenTwentyFour}{18.1}
\newcommand{\GflopsBlockedOneNTenTwentyFour}{3.52}
\newcommand{\GflopsIkjOneNTenTwentyFour}{3.95}
\newcommand{\SIkjSixNTwentyFortyEight}{5.17}
\newcommand{\SIkjTwelveNTwentyFortyEight}{5.51}
\newcommand{\SIkjBestNTwentyFortyEight}{5.51}
\newcommand{\PIkjBestNTwentyFortyEight}{12}
\newcommand{\GflopsIkjBestNTwentyFortyEight}{16.1}
\newcommand{\TseqIkjNTwentyFortyEight}{5.87}
\newcommand{\GflopsBlockedBestNTwentyFortyEight}{17.1}
\newcommand{\GflopsBlockedOneNTwentyFortyEight}{3.32}
\newcommand{\GflopsIkjOneNTwentyFortyEight}{3.26}
\newcommand{\RunOneScalarRatioNTenTwentyFour}{0.994}
\newcommand{\RunOneScalarRatioLoNTenTwentyFour}{0.986}
\newcommand{\RunOneScalarRatioHiNTenTwentyFour}{1.048}
\newcommand{\RunOneScalarRatioNTwentyFortyEight}{0.976}
\newcommand{\RunOneScalarRatioLoNTwentyFortyEight}{0.904}
\newcommand{\RunOneScalarRatioHiNTwentyFortyEight}{1.056}
\newcommand{\RunOneOverheadMinPct}{+19}
\newcommand{\RunOneOverheadMaxPct}{+26}
\newcommand{\RunOneIkjMinPct}{-16}
\newcommand{\RunOneIkjMaxPct}{-8}
\newcommand{\RunOneSLargeSix}{4.03}
\newcommand{\OverheadIkjMinPct}{-15.6}
\newcommand{\OverheadIkjMaxPct}{-10.1}
\newcommand{\SRelSixNTwoFiftySix}{5.26}
\newcommand{\SRelTwelveNTwoFiftySix}{5.43}
\newcommand{\SRelSixNFiveTwelve}{5.03}
\newcommand{\SRelTwelveNFiveTwelve}{5.19}
\newcommand{\SRelSixNTenTwentyFour}{4.69}
\newcommand{\SRelTwelveNTenTwentyFour}{4.90}
\newcommand{\SRelSixNTwentyFortyEight}{4.82}
\newcommand{\SRelTwelveNTwentyFortyEight}{5.14}
\newcommand{\SRelSixNThirtySeventyTwo}{4.82}
\newcommand{\SRelTwelveNThirtySeventyTwo}{4.98}
\newcommand{\OverheadMinPct}{+20.7}
\newcommand{\OverheadMaxPct}{+26.9}
\newcommand{\OverheadLarge}{1.214}
\newcommand{\OverheadLargeLo}{1.174}
\newcommand{\OverheadLargeHi}{1.234}
\newcommand{\VecAvxPOneNTenTwentyFour}{2.21}
\newcommand{\VecLoAvxPOneNTenTwentyFour}{2.13}
\newcommand{\VecHiAvxPOneNTenTwentyFour}{2.25}
\newcommand{\VecGflopsAvxPOneNTenTwentyFour}{9.3}
\newcommand{\VecFmaPOneNTenTwentyFour}{2.33}
\newcommand{\VecLoFmaPOneNTenTwentyFour}{2.27}
\newcommand{\VecHiFmaPOneNTenTwentyFour}{2.38}
\newcommand{\VecGflopsFmaPOneNTenTwentyFour}{9.7}
\newcommand{\VecScalarPOneNTenTwentyFour}{1.00}
\newcommand{\VecLoScalarPOneNTenTwentyFour}{0.97}
\newcommand{\VecHiScalarPOneNTenTwentyFour}{1.03}
\newcommand{\VecGflopsScalarPOneNTenTwentyFour}{4.2}
\newcommand{\VecAvxPSixNTenTwentyFour}{1.96}
\newcommand{\VecLoAvxPSixNTenTwentyFour}{1.87}
\newcommand{\VecHiAvxPSixNTenTwentyFour}{2.15}
\newcommand{\VecGflopsAvxPSixNTenTwentyFour}{39.5}
\newcommand{\VecThrAvxNTenTwentyFour}{4.26}
\newcommand{\VecFmaPSixNTenTwentyFour}{2.10}
\newcommand{\VecLoFmaPSixNTenTwentyFour}{1.94}
\newcommand{\VecHiFmaPSixNTenTwentyFour}{2.32}
\newcommand{\VecGflopsFmaPSixNTenTwentyFour}{42.3}
\newcommand{\VecThrFmaNTenTwentyFour}{4.35}
\newcommand{\VecScalarPSixNTenTwentyFour}{1.00}
\newcommand{\VecLoScalarPSixNTenTwentyFour}{0.89}
\newcommand{\VecHiScalarPSixNTenTwentyFour}{1.12}
\newcommand{\VecGflopsScalarPSixNTenTwentyFour}{20.1}
\newcommand{\VecThrScalarNTenTwentyFour}{4.81}
\newcommand{\VecAvxPOneNTwentyFortyEight}{2.15}
\newcommand{\VecLoAvxPOneNTwentyFortyEight}{2.07}
\newcommand{\VecHiAvxPOneNTwentyFortyEight}{2.31}
\newcommand{\VecGflopsAvxPOneNTwentyFortyEight}{9.0}
\newcommand{\VecFmaPOneNTwentyFortyEight}{2.30}
\newcommand{\VecLoFmaPOneNTwentyFortyEight}{2.07}
\newcommand{\VecHiFmaPOneNTwentyFortyEight}{2.43}
\newcommand{\VecGflopsFmaPOneNTwentyFortyEight}{9.6}
\newcommand{\VecScalarPOneNTwentyFortyEight}{1.00}
\newcommand{\VecLoScalarPOneNTwentyFortyEight}{0.95}
\newcommand{\VecHiScalarPOneNTwentyFortyEight}{1.06}
\newcommand{\VecGflopsScalarPOneNTwentyFortyEight}{4.2}
\newcommand{\VecAvxPSixNTwentyFortyEight}{1.90}
\newcommand{\VecLoAvxPSixNTwentyFortyEight}{1.54}
\newcommand{\VecHiAvxPSixNTwentyFortyEight}{2.58}
\newcommand{\VecGflopsAvxPSixNTwentyFortyEight}{39.1}
\newcommand{\VecThrAvxNTwentyFortyEight}{4.33}
\newcommand{\VecFmaPSixNTwentyFortyEight}{1.94}
\newcommand{\VecLoFmaPSixNTwentyFortyEight}{1.66}
\newcommand{\VecHiFmaPSixNTwentyFortyEight}{2.63}
\newcommand{\VecGflopsFmaPSixNTwentyFortyEight}{39.9}
\newcommand{\VecThrFmaNTwentyFortyEight}{4.14}
\newcommand{\VecScalarPSixNTwentyFortyEight}{1.00}
\newcommand{\VecLoScalarPSixNTwentyFortyEight}{0.73}
\newcommand{\VecHiScalarPSixNTwentyFortyEight}{1.36}
\newcommand{\VecGflopsScalarPSixNTwentyFortyEight}{20.5}
\newcommand{\VecThrScalarNTwentyFortyEight}{4.89}
\newcommand{\FreqOne}{3900}
\newcommand{\FreqTwo}{3700}
\newcommand{\FreqThree}{3400}
\newcommand{\FreqFour}{3300}
\newcommand{\FreqSix}{3150}
\newcommand{\FreqEight}{3150}
\newcommand{\FreqTwelve}{3050}
\newcommand{\FreqSeq}{3810}
\newcommand{\FreqRatioOne}{1.00}  % f(p)/f(1)
\newcommand{\FreqBoundOne}{1.00}  % min(p,6) * f(p)/f(1): speedup bound if only the frequency drop limited scaling
\newcommand{\FreqRatioTwo}{0.95}  % f(p)/f(1)
\newcommand{\FreqBoundTwo}{1.90}  % min(p,6) * f(p)/f(1): speedup bound if only the frequency drop limited scaling
\newcommand{\FreqRatioThree}{0.87}  % f(p)/f(1)
\newcommand{\FreqBoundThree}{2.62}  % min(p,6) * f(p)/f(1): speedup bound if only the frequency drop limited scaling
\newcommand{\FreqRatioFour}{0.85}  % f(p)/f(1)
\newcommand{\FreqBoundFour}{3.38}  % min(p,6) * f(p)/f(1): speedup bound if only the frequency drop limited scaling
\newcommand{\FreqRatioSix}{0.81}  % f(p)/f(1)
\newcommand{\FreqBoundSix}{4.85}  % min(p,6) * f(p)/f(1): speedup bound if only the frequency drop limited scaling
\newcommand{\FreqRatioEight}{0.81}  % f(p)/f(1)
\newcommand{\FreqBoundEight}{4.85}  % min(p,6) * f(p)/f(1): speedup bound if only the frequency drop limited scaling
\newcommand{\FreqRatioTwelve}{0.78}  % f(p)/f(1)
\newcommand{\FreqBoundTwelve}{4.69}  % min(p,6) * f(p)/f(1): speedup bound if only the frequency drop limited scaling
\newcommand{\CpuUtilOne}{1.00}
\newcommand{\CpuUtilTwo}{2.00}
\newcommand{\CpuUtilThree}{2.98}
\newcommand{\CpuUtilFour}{3.99}
\newcommand{\CpuUtilFive}{4.96}
\newcommand{\CpuUtilSix}{5.94}
\newcommand{\CpuUtilEight}{6.32}
\newcommand{\CpuUtilTen}{9.05}
\newcommand{\CpuUtilTwelve}{11.88}
\newcommand{\CpuUtilSixteen}{10.41}
\newcommand{\CpuUtilTwentyFour}{11.66}
\newcommand{\RssSeqNTwoFiftySix}{3.0}
\newcommand{\RssMaxNTwoFiftySix}{3.5}
\newcommand{\ArrayMibNTwoFiftySix}{1.5}
\newcommand{\RssSeqNFiveTwelve}{7.5}
\newcommand{\RssMaxNFiveTwelve}{8.1}
\newcommand{\ArrayMibNFiveTwelve}{6.0}
\newcommand{\RssSeqNTenTwentyFour}{25.5}
\newcommand{\RssMaxNTenTwentyFour}{25.9}
\newcommand{\ArrayMibNTenTwentyFour}{24.0}
\newcommand{\RssSeqNTwentyFortyEight}{97.5}
\newcommand{\RssMaxNTwentyFortyEight}{98.0}
\newcommand{\ArrayMibNTwentyFortyEight}{96.0}
\newcommand{\RssSeqNThirtySeventyTwo}{217.5}
\newcommand{\RssMaxNThirtySeventyTwo}{218.0}
\newcommand{\ArrayMibNThirtySeventyTwo}{216.0}


% Draft helper: remove all \todo{} before submission (report-review checks this)
\newcommand{\todo}[1]{\textcolor{red}{[TODO: #1]}}

\begin{document}

\title{Thread-Level Scaling of Shared-Memory Matrix Multiplication on a Six-Core Processor\\[0.3em]
\large Big Data -- Individual Assignment 03}

\author{\IEEEauthorblockN{Daniel Rodríguez Alonso}
\IEEEauthorblockA{Grado en Ciencia e Ingeniería de Datos\\
Universidad de Las Palmas de Gran Canaria (ULPGC)\\
Big Data, Academic Year 2026--2027}}

\maketitle

% ------------------------------------------------------------------ Abstract
\begin{abstract}
ABSTRACT
\end{abstract}

\begin{IEEEkeywords}
matrix multiplication, OpenMP, shared-memory parallelism, speedup, parallel efficiency,
simultaneous multithreading, SIMD vectorisation, benchmarking
\end{IEEEkeywords}

% ------------------------------------------------------------------ Introduction
\section{Introduction and Objective}
\label{sec:intro}
Dense matrix multiplication is the reference kernel of numerical and data-intensive computing: it
performs $2n^3$ floating-point operations on $3n^2$ values, so it has enough work per byte to keep
every core of a processor busy. Assignments~1 and~2 of this course measured single-threaded
implementations; the cache-blocked kernel of Assignment~2 reached its best performance at a
tile size of 256 on the same laptop used here. A modern laptop, however, has several cores and two
hardware threads per core, and a single-threaded kernel uses a small fraction of that capacity.
This study measures how much of it a straightforward shared-memory parallelisation recovers.

\textbf{Research question (RQ).} \emph{On a six-core, twelve-thread laptop processor
(Intel Core i7-9750H), how does the number of OpenMP threads $p$ affect the runtime, speedup and
efficiency of a cache-blocked dense matrix multiplication relative to an equivalent sequential
baseline, and for which matrix sizes $n$ does parallel execution pay off?} The factor varied is
$p\in\{1,2,3,4,5,6,8,10,12,16,24\}$, with $n\in\{256,\dots,3072\}$ as a second factor; the outcomes
are the kernel runtime, its variability, the speedup $S(p)$, the efficiency $E(p)$ and the peak memory.
Two secondary experiments vary the memory behaviour of the kernel (blocked vs.\ un-tiled) and the
SIMD code generation at a fixed thread count (optional vectorisation extension).

The hypotheses below were written before the main benchmark (\texttt{docs/experiment\_plan.md}) and
are predictions, not results:
\begin{itemize}
\item[\textbf{H1}] For large $n$, $S(p)$ grows almost linearly up to the six physical cores but stays
  below $p$, partly because the all-core turbo frequency is lower than the single-core one.
\item[\textbf{H2}] The second hardware thread of each core (simultaneous multithreading, SMT;
  $7\le p\le 12$) adds little speedup for this compute-bound kernel.
\item[\textbf{H3}] Oversubscription ($p=16,24$, more threads than the 12 logical CPUs) gives no gain
  and may be slower.
\item[\textbf{H4}] For small $n$ (256) fork/join overhead limits $S(p)$ well below the large-$n$ values.
\item[\textbf{H5}] The parallel binary with one thread is within a few percent of the sequential baseline.
\item[\textbf{H6}] The un-tiled $ikj$ kernel scales worse than the blocked kernel at $n=2048$, where one
  operand (32\,MiB) exceeds the 12\,MiB last-level cache.
\item[\textbf{H7}] AVX2 vectorisation gives a further speedup at fixed $p$, with a smaller relative gain at
  $p=6$ than at $p=1$ if the vectorised kernel becomes limited by memory bandwidth.
\end{itemize}

% ------------------------------------------------------------------ Background
\section{Background and Related Work}
\label{sec:background}
\textbf{Kernel.} The classical algorithm computes $C_{ij}=\sum_k A_{ik}B_{kj}$ with $2n^3$
floating-point operations. With row-major storage, the $i$-$k$-$j$ loop order makes the innermost loop
a unit-stride update $c_{i,\cdot}\mathrel{+}=a_{ik}\,b_{k,\cdot}$; loop tiling (blocking) additionally
reuses a tile of $B$ from cache across many rows, reducing the traffic to main memory once the
operands no longer fit in cache~\cite{lam1991}. High-performance libraries combine blocking for every
cache level with hand-written SIMD micro-kernels~\cite{goto2008}; this study deliberately stays with a
portable C kernel and the compiler's code generation.

\textbf{Shared-memory parallelism.} OpenMP~\cite{dagum1998,openmp52} parallelises a C program with
compiler directives: a \texttt{parallel} region forks a team of threads that share the address space,
and an implicit barrier joins them at the end. The thread count is set with \texttt{OMP\_NUM\_THREADS};
\texttt{OMP\_PLACES} and \texttt{OMP\_PROC\_BIND} control the CPUs the threads are bound
to~\cite{openmp52}. Partitioning the rows of $C$ among threads gives each thread disjoint output rows,
so no locks or reductions are needed.

\textbf{Metrics and limits.} With $T_{seq}(n)$ the time of the sequential baseline and $T_p(n)$ that of
the parallel version with $p$ threads, $S(p)=T_{seq}/T_p$ and $E(p)=S(p)/p$. Amdahl's
law~\cite{amdahl1967} bounds $S$ by the serial fraction of the work; Gustafson~\cite{gustafson1988}
noted that the parallel fraction usually grows with the problem size, which is why larger $n$ is
expected to scale better. The roofline model~\cite{williams2009} relates attainable performance to the
arithmetic intensity of a kernel and the memory bandwidth: a kernel with poor cache reuse saturates the
shared memory bus before all cores are busy. SMT (Intel Hyper-Threading) exposes two logical
processors per core that share the core's execution units and caches~\cite{marr2002}, so it helps
mostly when one thread alone leaves those units idle. On this processor the maximum turbo frequency
(4.5\,GHz) is far above the base frequency (2.6\,GHz)~\cite{intel9750h}, and the frequency actually
reached depends on how many cores are active, which is a known confounder of speedup measurements.

\textbf{Benchmarking.} Hoefler and Belli~\cite{hoefler2015} recommend reporting the summary statistic,
its variability and the exact configuration of every parallel measurement, and computing speedups
against a well-tuned sequential baseline rather than against the parallel code with one thread.
Mytkowicz et al.~\cite{mytkowicz2009} show that seemingly innocuous setup details bias measurements,
which motivates the interleaved and randomised run order used here. Floating-point error bounds for
inner products follow Higham~\cite{higham2002}.

% ------------------------------------------------------------------ Methodology
\section{Methodology}
\label{sec:method}

\subsection{Implementations}
All kernels are written in C11 in one source file, \texttt{src/c/matmul.c}, for float64 matrices
stored row-major and contiguous, $C=AB$ overwriting $C$. No external numerical library is linked.
\begin{itemize}
\item \textbf{Blocked kernel} (the kernel studied): loops $ii\!\to\!kk\!\to\!jj$ over $b\times b$ tiles
  with an $i\!\to\!k\!\to\!j$ tile update inside; the tile size $b=256$ was selected by a pilot
  (\cref{sec:protocol}). For every $C_{ij}$ the products are accumulated in increasing $k$, as in the
  plain $ijk$ loop.
\item \textbf{Un-tiled $ikj$ kernel}: the same innermost update without tiling (secondary experiment).
\item \textbf{Naive $ijk$ kernel}: textbook loop, used only as the trusted reference in the tests.
\end{itemize}
\textbf{Equivalent sequential baseline.} The baseline is the \emph{same} source compiled without
\texttt{-fopenmp} (\texttt{bin/matmul\_seq}): the OpenMP directives are ignored and the team-size
helpers return a one-thread team, so algorithm, loop order, tile size, layout, numeric type, compiler
flags (\texttt{-O2}) and inputs are identical; only the OpenMP runtime differs. The sequential binary
does not link \texttt{libgomp} (\texttt{results/evidence/linkage.txt}). For consistency with
Assignment~2, its single-threaded blocked kernel uses the same loop nest and tile size.

\textbf{Parallel version} (\texttt{bin/matmul\_omp}, \texttt{-O2 -fopenmp}). One
\texttt{omp parallel} region per call; thread $t$ of a team of $p$ owns the contiguous rows
$[\lfloor nt/p\rfloor,\lfloor n(t+1)/p\rfloor)$ of $C$, zeroes them and runs the blocked loops over
them. The partition balances the work to within one row, writes are disjoint, $A$ and $B$ are only
read, and the implicit barrier at the end of the region guarantees that $C$ is complete on return.
The $ikj$ kernel uses \texttt{omp parallel for schedule(static)} over rows, which gives the same
contiguous partition. Because each element is computed by one thread with the same sequence of
operations, the parallel result is bit-identical to the sequential one for every $p$.

\textbf{Vectorisation variants} (same source, OpenMP builds, only flags differ): \emph{scalar}
\texttt{-O2 -fno-tree-vectorize}; \emph{AVX2} \texttt{-O3 -march=native}; \emph{AVX2+FMA}
\texttt{-O3 -march=native -ffp-contract=fast}. All builds use \texttt{-std=c11}; in this ISO mode gcc
does not contract $a\cdot b+c$ into fused multiply-adds (FMA) unless \texttt{-ffp-contract=fast} is
given~\cite{gccopt13}, so the AVX2 build performs the same rounding as the scalar one and FMA is
studied as a separate variant.

\subsection{Inputs}
\texttt{scripts/gen\_inputs.py} draws $A$ and then $B$ from $\mathcal{U}[-1,1)$ with NumPy's PCG64
generator (\texttt{default\_rng(42)}) in float64 and stores them as raw little-endian files with
SHA-256 checksums (\texttt{data/inputs/checksums.sha256}). Every method, thread count and repetition
of a given $n$ reads the same two files; loading happens outside the timed region.

\subsection{Experimental environment}
\Cref{tab:env} summarises the machine (full capture in \texttt{results/environment.md}). Core $k$
corresponds to logical CPUs $\{k,k+6\}$. With \texttt{OMP\_PLACES=cores} and
\texttt{OMP\_PROC\_BIND=close}, thread $t$ is bound to core $t\bmod 6$ (both of its logical CPUs):
up to $p=6$ every thread has a core of its own, for $7\le p\le12$ two threads share a core (SMT), and
for $p>12$ the logical CPUs are \emph{oversubscribed}. The bindings were verified with
\texttt{OMP\_DISPLAY\_AFFINITY} for every $p$ (\texttt{results/evidence/affinity\_display.txt}) and the
real team size and the CPU of each thread are recorded in every measurement; in all
\NumRowsMeasured{} measured rows the team size equals the requested $p$ (\NumThreadMismatch{}
mismatches). The sequential binary is pinned with \texttt{taskset -c 0,6}, i.e.\ to the same physical
core as OpenMP place~0, so $T_{par}(1)/T_{seq}$ isolates the runtime overhead. Internal library
threads are disabled (\texttt{OPENBLAS\_NUM\_THREADS=MKL\_NUM\_THREADS=1}; no BLAS is linked) and
\texttt{OMP\_DYNAMIC=false}.

\begin{table}[t]
\centering
\caption{Execution environment (captured by \texttt{capture\_env.sh} on 2026-11-13).}
\label{tab:env}
\footnotesize
\setlength{\tabcolsep}{3pt}
\begin{tabular}{@{}lp{0.72\columnwidth}@{}}
\toprule
Item & Value \\
\midrule
CPU & Intel Core i7-9750H, 1 socket, 6 cores / 12 logical CPUs (SMT), base 2.6\,GHz, max.\ turbo
      4.5\,GHz~\cite{intel9750h}; AVX2, FMA \\
Caches & L1d 32\,KiB and L2 256\,KiB per core; L3 12\,MiB shared \\
Memory & 15.5\,GiB DDR4, 1 NUMA node \\
OS & Ubuntu 24.04.5 LTS, Linux 7.0.0-42; \texttt{intel\_pstate} governor \texttt{powersave},
     turbo enabled, laptop on AC power \\
Compiler & gcc 13.3.0, \texttt{-std=c11}; OpenMP runtime libgomp (gcc 13) \\
Analysis & Python 3.12.3, NumPy 2.5.4, pandas 3.1.0, Matplotlib 3.11.2 \\
Threads & \texttt{OMP\_NUM\_THREADS}=$p$, \texttt{OMP\_PLACES=cores}, \texttt{OMP\_PROC\_BIND=close},
          \texttt{OMP\_DYNAMIC=false}; seq.\ pinned to CPUs 0,6 \\
\bottomrule
\end{tabular}
\end{table}

\subsection{Correctness validation}
\texttt{make test} runs, for every build variant, both kernels on identity and zero matrices,
hand-checked $2\times2$ and $3\times3$ integer products (exact in float64), odd sizes, $n<p$,
$n$ not a multiple of the tile, tile sizes 1, 7, 64, 200 and 256, and random matrices against the naive
$ijk$ reference. Random results are accepted if $|C-\hat C|\le 2\gamma_n|A||B|$ componentwise, with
$\gamma_n=nu/(1-nu)$ and $u=2^{-53}$; this is twice the classical worst-case bound for inner products of
length $n$~\cite{higham2002}, which covers the two computations being compared. The test also checks that
the parallel result is bit-identical to the sequential one for $p\in\{1,2,6,7,12,24\}$, and an
end-to-end script compares every binary against NumPy. All checks pass
(\texttt{results/logs/make\_test.log}); the largest observed error is $6.4\times10^{-3}$ of the bound
(FMA build: $4.0\times10^{-3}$, not bit-identical, as expected from the different rounding). During the
benchmark the checksum of $C$ is recorded after every repetition: across all scalar and AVX2 builds and
all $p$ there is \ChecksumDistinctMax{} distinct checksum per $(n,\text{kernel})$.

\textbf{Safe access to shared data.} Besides the disjoint-row argument, a ThreadSanitizer build of the
harness was run for both kernels with $p=2,6,12$ (\texttt{results/evidence/tsan.txt}). It reports
warnings, but all of them are on stack addresses (the block of shared-variable pointers that libgomp
passes to the threads, and the team-size probe), and none is inside the tile update that writes $C$. The
system libgomp is not instrumented, so TSan cannot see its barriers; these warnings are the known false
positives of that setting and are not evidence of a race in the kernel.

\subsection{Measurement protocol and budgets}
\label{sec:protocol}
\textbf{Timer and boundary.} \texttt{clock\_gettime(CLOCK\_MONOTONIC)} around exactly one kernel call
(\emph{kernel} boundary): zeroing $C$, the product and, for OpenMP, the fork/join of that call's team.
The thread team is created before timing by a probe region, $C$ is touched before timing, and the
checksum of $C$ is computed after each repetition (outside the timer) so that the result is consumed.

\textbf{Repetitions and order.} Every process run goes through \texttt{run\_measured.py}, which takes
a machine-wide lock (other agents share the laptop), records the peak resident set size (RSS) with
GNU \texttt{time} and enforces the 600\,s time-out. Each configuration is executed in $R=5$ process
runs; inside each run $w$ untimed warm-ups are followed by $m$ measured repetitions, with
$(w,m)=(1,30),(1,10),(1,3),(0,2),(0,1)$ for $n=256,512,1024,2048,3072$, i.e.\ $k=150,50,15,10,5$
measured repetitions per configuration (short kernels get more repetitions to reduce timer noise;
long ones need no warm-up because the inputs are already in memory). The full list of configurations
is run five times, shuffled differently in each round, so that slow drifts (temperature, background
activity) are spread over all configurations~\cite{mytkowicz2009}.

\textbf{Statistics.} $T_{seq}$ and $T_p$ are the \emph{medians} of the measured repetitions; variability
is the interquartile range (IQR). $S(p)$ is the ratio of medians with a 95\% bootstrap confidence
interval (CI; 2000 resamples of both samples); $E(p)=S(p)/p$. Memory is the peak RSS of each process
(GNU \texttt{time} \%M), compared with the analytical array storage $3\cdot8n^2$ bytes.

\textbf{Interference.} The harness records, for every repetition, the CPU time consumed by \emph{other}
processes divided by the wall time (\texttt{other\_busy\_cpus}, from \texttt{/proc/stat}). All rows are
kept; a sensitivity analysis recomputes $S(p)$ after removing repetitions longer than 0.5\,s with more
than \ContaminationThreshold{} busy foreign CPU.

\textbf{Frequency evidence.} Hardware performance counters are not accessible on this machine
(\texttt{perf\_event\_paranoid}=4). A separate \texttt{freq\_probe} experiment ($n=2048$, 2 rounds) samples
the kernel's cpufreq estimate of the effective frequency of the busy CPUs every 0.1\,s; these runs are
not used for timing.

\textbf{Tile size (pilot).} The pilot compared $b\in\{32,64,128,256\}$ for the sequential kernel and
$b\in\{64,128,256\}$ for $p=6$ at $n=512$--$2048$ (\texttt{results/raw/pilot\_tile*.csv}). $b=64$ was
fastest at $n\le1024$ by at most 5--8\%, while $b=256$ was clearly fastest at $n=2048$ (sequential
4.39\,s vs.\ 5.49\,s with $b=64$). $b=256$ (the smallest worst-case loss, and the value selected in
Assignment~2) was fixed for all experiments, so that the sequential baseline is not handicapped at
large $n$, which would inflate the speedups.

\textbf{Sizes and budgets.} Main sweep: blocked kernel, $n\in\{256,512,1024,2048,3072\}$, all 11 thread
counts plus the baseline. $ikj$: $n\in\{1024,2048\}$, same thread counts. Vectorisation:
$n\in\{1024,2048\}$, $p\in\{1,6\}$. Budgets agreed before the runs: at most 600\,s per process,
8\,GB of memory ($3\cdot8\cdot3072^2$\,B $=216$\,MiB is far below) and 90\,min of machine time.
A configuration that fails is recorded with its status and never assigned a time.

% ------------------------------------------------------------------ Results
\section{Results}
\label{sec:results}
RESULTS

% ------------------------------------------------------------------ Discussion
\section{Discussion}
\label{sec:discussion}
DISCUSSION

% ------------------------------------------------------------------ Conclusions
\section{Conclusions}
\label{sec:conclusions}
CONCLUSIONS

% ------------------------------------------------------------------ Materials + AI statement
\section*{Reproducibility Materials}
Repository: \url{https://github.com/danirodriguezz/bigdata-assignment03-parallel-matmul}
(to be pushed by the author; the submitted version is the last commit of \texttt{main}). The archive
\texttt{Assignment03\_RodriguezAlonso\_Materials.zip} contains the C sources, tests, \texttt{Makefile},
configuration (\texttt{configs/experiment.json}), input generator and checksums, the experiment
driver and \texttt{run\_measured.py}, all raw CSVs (\texttt{results/raw/}), logs, evidence files
(vectorisation report, disassembly, affinity, TSan), the analysis script and this report's sources.
\texttt{make test inputs bench} reproduces the measurements (about one hour on this machine) and
\texttt{make analysis report} regenerates every table, figure and number from the raw data.

\section*{Use of Generative AI}
AISTATEMENT

% ------------------------------------------------------------------ References
\bibliographystyle{IEEEtran}
\bibliography{references}

\end{document}
